\fontfam[LM]
\uselanguage{czech}
\def\begcitation{\par\medskip\leftskip=2em\rightskip=2em\noindent}
\def\endcitation{\par\leftskip=0em\rightskip=0em\medskip}
% odborny — academic / non-fiction publication
% Termes (Times-like), A4, wider outer margin for margin notes, folio in footer.
\fontfam[Termes]
\typosize[11/13]
\margins/2 a4 (30,45,25,25)mm
\footline={\hss\folio\hss}



\chap Ukázkový dokument
Toto je ukázkový Markdown soubor pro testování nástroje md2opmac.
Obsahuje příklady všech podporovaných konstrukcí.


\sec Formátování textu
Tento odstavec obsahuje {\bf tučný text} a~{\it kurzívu}. Lze také kombinovat:
{\bf tučná {\it tučná kurzíva} tučná}. Inline kód vypadá takto: {\tt int main()}.

Druhý odstavec v~tomto oddíle. Text může obsahovat {\tt proměnné} nebo {\tt funkce()}.


\sec České typografické konvence
Nástroj automaticky aplikuje české typografické pravidla.

Uvozovky: \uv{toto je v~uvozovkách} a~také \uv{toto je již česky}.

Uvozovky se~seznamem a~pomlčkou:

\begitems
* \uv{Poslouchej svůj hlas~--- vypadá to na odpor, ale je v~něm touha.}
* \uv{Druhá položka s~pomlčkou~--- taky funguje.}
\enditems

Pomlčka odděluje části věty~-- například takto. Nebo pomocí em-dash~--- takhle.

Jednopísmenné předložky: v~lese, z~kopce, s~kamarádem, k~domovu, u~řeky.
Spojky: Jan a~Marie, Petr i~Pavel, slunce o~půlnoci.

Výpustka na konci věty\dots{}


\secc Podpodnadpis s~příkladem
Text na úrovni třetí nadpisové úrovně. Zde je odkaz na
\ulink[https://petr.olsak.net/optex/]{stránku projektu OpTeX} a~ještě jeden
\ulink[https://example.com]{odkaz s~titulkem}.


\sec Seznamy
Nečíslovaný seznam:

\begitems
* první položka
* druhá položka s~{\bf tučným} textem
* třetí položka s~{\tt kódem}
\enditems

Číslovaný seznam:

\begitems \style n
* první krok
* druhý krok
* třetí krok
\enditems


\sec Blok kódu
Příklad Rust kódu:

\begtt
fn main() {
    println!("Ahoj, světe!");
    let x = 42;
    let y = x * 2;
    println!("Výsledek: {}", y);
}
\endtt

Příklad bez specifikace jazyka:

\begtt
$ cargo build --release
$ ./target/release/md2opmac vstup.md -o vystup.tex
\endtt


\sec Citace
\begcitation
\uv{Filozofie je příprava na smrt,} napsal Platón.
Přijmi ji, i~když se~ti zdá teoretická.

\endcitation

\begcitation
\uv{Kdo se~bojí, nesmí do lesa,} říká staré přísloví.
Ale \uv{odvaha není absence strachu~-- je to rozhodnutí, že něco jiného je důležitější než strach.}

\endcitation

\begcitation
Druhý odstavec citace obsahuje {\it formátování} a~{\tt kód}.

\endcitation


\sec Horizontální linka
Text před oddělovačem.

\noindent\hrule

Text po oddělovači.


\sec Tabulka
Ukázka GFM tabulky se~zarovnáním sloupců:

\table{lcr}{
Název & Typ & Hodnota \crli
alfa & string & 1 \cr
beta & integer & 42 \cr
gama & boolean & 0 \cr
}


\sec Obrázek
Pokud obrázek existuje, změří se~a vloží automaticky:

\begtt
![Popisek](cesta/k/obrazku.png)
\endtt

Pokud soubor neexistuje, vypíše se~varování na stderr a~použije se~{\tt \char92 picw=\char92 hsize}.


\bye
